% INTERNAL DRAFT -- NOT FOR DISTRIBUTION, NOT TO BE POSTED OR UPLOADED.
% Computations referenced as computations/{homology,sage,milp}/... (repository-relative).
\documentclass[11pt]{amsart}
\usepackage[T1]{fontenc}
\usepackage{lmodern}
\usepackage{amsmath,amssymb,amsthm,mathtools}
\usepackage[expansion=false]{microtype}
\usepackage{booktabs,array}
\usepackage{xcolor}
\usepackage[a4paper,margin=28mm]{geometry}
\usepackage[backend=biber,style=alphabetic,maxbibnames=9,sorting=nyt]{biblatex}
\addbibresource{refs.bib}
\usepackage{hyperref}
\hypersetup{colorlinks=true,linkcolor=blue!45!black,citecolor=blue!45!black,urlcolor=blue!45!black}
\setlength{\emergencystretch}{2em}

\newcommand{\todo}[1]{{\color{red}\textbf{[TODO:} #1\textbf{]}}}

\numberwithin{equation}{section}
\newtheorem{theorem}{Theorem}[section]
\newtheorem{lemma}[theorem]{Lemma}
\newtheorem{proposition}[theorem]{Proposition}
\newtheorem{corollary}[theorem]{Corollary}
\newtheorem{mainthm}{Theorem}
\renewcommand{\themainthm}{\Alph{mainthm}}
\theoremstyle{definition}
\newtheorem{definition}[theorem]{Definition}
\newtheorem{problem}[theorem]{Problem}
\theoremstyle{remark}
\newtheorem{remark}[theorem]{Remark}

\newcommand{\CP}{\mathbf{CP}}
\newcommand{\Z}{\mathbf Z}
\newcommand{\Q}{\mathbf Q}
\newcommand{\R}{\mathbf R}
\newcommand{\C}{\mathbf C}
\newcommand{\F}{\mathbf F}
\newcommand{\PP}{\mathbf P}
\DeclareMathOperator{\Tor}{Tor}
\DeclareMathOperator{\SM}{SM}
\DeclareMathOperator{\Morse}{Morse}
\DeclareMathOperator{\Crit}{Crit}
\DeclareMathOperator{\Fix}{Fix}
\DeclareMathOperator{\rank}{rank}
\DeclareMathOperator{\Bl}{Bl}
\DeclareMathOperator{\Sym}{Sym}
\DeclareMathOperator{\Gr}{Gr}
\newcommand{\id}{\mathrm{id}}

\title[Fixed points below the stable Morse number in dimension ten]{Hamiltonian fixed points below the stable Morse number in dimension ten, and an obstruction in dimension six}
\author{[AUTHOR NAME]}
\date{October 9, 2026 (internal; not for distribution)}

\begin{document}

\begin{abstract}
We construct a simply connected closed K\"ahler manifold $M$ of real
dimension ten and a Hamiltonian diffeomorphism of $M$ with only
nondegenerate fixed points whose number is $\lambda(M)-4=\SM(M)-4$, where
$\lambda$ is the integral Morse count and $\SM$ is the stable Morse number.
The count equals the cyclic integral Floer bound of Bai--Xu, which is
therefore sharp already in complex dimension five. The new ingredients are
a simply connected threefold $F$ with Brauer-type $2$-torsion, obtained as a
resolved $\Z/5$-quotient of an invariant Artin--Mumford double solid, a
divisor $S\subset F$ with $\pi_1(S)=\Z/5$, and a blow-up along a center
that meets the fixed locus of the Hamiltonian involution transversally, so
that no new fixed components appear. In the other direction, the Bai--Xu bound excludes simply connected
counterexamples in real dimension six, and in real dimension eight every
construction of the type of Proposition~\ref{prop:perturb} (a finite-order
element of a Hamiltonian torus action plus a small invariant perturbation)
yields at least $\lambda(M)$ fixed points. Among simply connected closed
symplectic manifolds of real dimension at least six, the minimal dimension
of a counterexample to the stable Morse bound is therefore $8$ or $10$.
\end{abstract}

\maketitle

\noindent\textbf{Status.} Internal draft. Theorem numbers cited from
\cite{OpenAI22} refer to its October~5, 2026 build. Every point that has
not been checked against the primary literature is marked \todo{\dots}.
Nothing in this document has been circulated.

\tableofcontents

%======================================================================
\section{Introduction}\label{sec:intro}

\subsection{Stable Morse number and Hamiltonian fixed points}
Let $W$ be a connected closed smooth manifold. Its \emph{stable Morse
number} $\SM(W)$ is the minimum of $\#\Crit(\Phi)$ over $k\ge 0$ and Morse
functions $\Phi\colon W\times\R^k\to\R$ that coincide with a nondegenerate
quadratic form in the $\R^k$-variables outside a compact set; the case
$k=0$ is the ordinary Morse number $\Morse(W)\ge\SM(W)$. For a closed
symplectic manifold $(M,\omega)$ and a one-periodic Hamiltonian $H$ all of
whose fixed points are nondegenerate, the \emph{stable Morse formulation}
of the Arnold problem asks whether
\begin{equation}\label{eq:SMconj}
  \#\Fix_0(\phi_H^1;H)\ \ge\ \SM(M),
\end{equation}
where $\Fix_0$ denotes fixed points with contractible orbits; see
\cite[Conjecture~1.2]{Golovko2020}. Positive results are known under
asphericity hypotheses \cite{DRG2017}. Counterexamples on simply connected
K\"ahler manifolds were found in \cite{OpenAIStable,OpenAI22,OpenAICompanion},
the smallest previously recorded fixed dimension being real dimension $22$
\cite{OpenAI22}; a non-simply-connected counterexample to the ordinary
Morse-number bound in real dimension $12$ is given in \cite{OpenAISept23}.

For a closed manifold $W$ put
\[
  b(W)=\sum_i\rank H_i(W;\Z),\qquad t_i(W)=d(\Tor H_i(W;\Z)),\qquad
  \lambda(W)=b(W)+2\sum_i t_i(W),
\]
where $d$ is the minimal number of generators. Then $\SM(W)\ge\lambda(W)$
for every closed $W$ \cite[Prop.~2.2]{OpenAI22}, with equality for simply
connected $W$ of dimension $\ge 6$ by Smale's theorem
(\S\ref{subsec:morse}).

\subsection{Main results}
The following theorem lowers the dimension of a simply connected
counterexample from $22$ to $10$.

\begin{mainthm}\label{thmA}
There exist a simply connected closed K\"ahler manifold $(M,\omega)$ of
complex dimension $5$ and a smooth one-periodic Hamiltonian $H$ on $M$
such that every fixed point of $\phi_H^1$ is nondegenerate and
\[
  \#\Fix(\phi_H^1)=\#\Fix_0(\phi_H^1;H)=\lambda(M)-4=\SM(M)-4=\Morse(M)-4 .
\]
Moreover the minimal Chern number of $M$ is $1$ and
$\#\Fix(\phi_H^1)$ equals the cyclic integral Floer bound
$\lambda_1(M)$ of Bai--Xu (Definition~\ref{def:lambdaN}); this bound is
therefore attained in complex dimension $5$.
\end{mainthm}

The gap $4$ does not depend on any free Betti number. For the resolution
chosen in \S\ref{sec:F} we obtain $\SM(M)=524$ and $\#\Fix=520$
(Remark~\ref{rem:absolute}).

\begin{mainthm}\label{thmB}
Let $(M,\omega)$ be a closed symplectic manifold with a Hamiltonian action
of a torus $T$, and let $g\in T$ have finite order. Suppose that every
connected component $C_1,\dots,C_r$ of $\Fix(g)$ has real dimension at
most $6$ (this holds whenever $\dim_\R M=8$ and $g$ does not act
trivially). Then
\[
  \sum_{j=1}^r\lambda(C_j)\ \ge\ \lambda(M).
\]
Consequently, if $\dim_\R M=8$, every Hamiltonian diffeomorphism
$\psi^\epsilon\circ g$ produced by Proposition~\ref{prop:perturb} has at
least $\lambda(M)$ nondegenerate fixed points, and at least $\SM(M)$ if
$M$ is simply connected.
\end{mainthm}

Theorem~\ref{thmB} is complemented by necessary conditions valid for every
simply connected eight-dimensional counterexample
(Proposition~\ref{prop:nec8}): the minimal Chern number must be $1$, and
the dominant primes of $\Tor H_2(M;\Z)$ and $\Tor H_3(M;\Z)$ must differ.
All previously known simply connected constructions
\cite{OpenAIStable,OpenAI22,OpenAICompanion} and the construction of
Theorem~\ref{thmA} are of the type excluded by Theorem~\ref{thmB} in
dimension $8$. In dimension $6$ the Bai--Xu bound already excludes
simply connected counterexamples (Proposition~\ref{prop:six}); real
dimensions $2$ and $4$ are discussed in \S\ref{sec:six}. In summary:

\begin{corollary}\label{cor:mindim}
Among simply connected closed symplectic manifolds of real dimension at
least $6$, the minimal dimension of a counterexample to
\eqref{eq:SMconj} is $8$ or $10$.
\end{corollary}

\begin{table}[h]
\centering\small
\begin{tabular}{@{}lll@{}}
\toprule
real dim & status & reference\\
\midrule
$2$ & no counterexample & \S\ref{sec:six}\\
$4$ & simply connected: open, reduces to $\SM$ vs.\ $\lambda$ & \S\ref{sec:six}\\
$6$ & no simply connected counterexample & Prop.~\ref{prop:six}\\
$8$ & open; excluded within the finite-order/torus mechanism & Theorem~\ref{thmB}, Prop.~\ref{prop:nec8}\\
$10$ & simply connected K\"ahler counterexample, Floer bound sharp & Theorem~\ref{thmA}\\
$12$ & non-simply-connected counterexample to the Morse-number bound & \cite{OpenAISept23}\\
$12$, $16$ & simply connected K\"ahler counterexamples & Remarks~\ref{rem:dim12}, \ref{rem:dim16}\\
$22$ & simply connected K\"ahler family, unbounded deficit & \cite{OpenAI22}\\
\bottomrule
\end{tabular}
\caption{Current status by dimension (nondegenerate Hamiltonians).}
\label{tab:status}
\end{table}

\subsection{Previous work versus our contributions}
\emph{Previous work.} The mechanism of \cite{OpenAIStable,OpenAI22,OpenAICompanion}
takes $M=X\times Y$, where $X$ is a simply connected K\"ahler manifold
with torsion at two primes and $Y$ is $\CP^1\times\CP^1$ blown up at four
torus-fixed points; the half-turn $g_Y$ of a Hamiltonian circle action on $Y$
fixes four lines, and a small invariant perturbation of $\id_X\times g_Y$
has $4\lambda(X\times\CP^1)$ fixed points. Because $d$ takes the
\emph{maximum} over primes in each degree, convolution of the torsion
profile of $X$ with the free profiles $(1,1)$ of $\CP^1$ and $(1,6,1)$ of
$Y$ can make $\lambda(X\times Y)>4\lambda(X\times\CP^1)$. The torsion of
$X$ comes from blow-ups of a free ambient space along bielliptic surfaces
and their products, and the resulting examples live in real dimension
$22$ \cite{OpenAI22} (and $3332$ in \cite{OpenAICompanion}, where the
cyclic Floer bound is attained). Independently, \cite[Thm.~1.1]{OpenAISept23}
gives a nondegenerate counterexample to the \emph{ordinary} Morse-number
bound on a closed symplectic manifold of real dimension $12$ with
fundamental group $A_5^k$ (not simply connected), with
$\#\Fix=\chi(M)+128<\Morse(M)$; there the effect comes from $\pi_1$, not
from torsion. A degenerate six-dimensional example on a quadric with three
fixed points, for the critical-point count, is in \cite{OpenAIDegenerate}.
The stable Morse bound is formulated in \cite{Golovko2020}; Bai--Xu
\cite{BaiXu2022,BaiXu2026} prove the cyclic integral lower bound. Ma's
preprint \cite[Thm.~1.5]{Ma2013} asserts that the minimal number of fixed
points equals the Morse number; see \S\ref{subsec:Ma}.

\emph{Where the previous method fails.} The product mechanism cannot work
below real dimension $12$ (Proposition~\ref{prop:productfail}): for a
simply connected $X$ of real dimension $\le 6$ the torsion of $X$ is a
single group placed in degrees $2,3$, and no exchange of dominant prime
can occur. Moreover blow-ups of a free ambient space never create torsion
in degree $2$, which is forced in dimension $8$ by
Proposition~\ref{prop:nec8}.

\emph{Our contributions.}
\begin{enumerate}
\item A new torsion source: a simply connected threefold $F$ with
  $\Tor H_2(F)=\Tor H_3(F)=\Z/2$ and \emph{no} $5$-torsion, containing a
  smooth divisor $S$ with $\pi_1(S)=\Z/5$ (\S\ref{sec:F}). The proof that
  the $5$-torsion vanishes uses a permutation-module structure of
  $H_2$ of an Artin--Mumford double solid and Lefschetz duality.
\item A non-product construction: the blow-up of $F\times Y$ along
  $W=S\times(L_0\sqcup L_2)$, where $L_0,L_2\subset Y$ are invariant lines
  meeting the fixed curves transversally. The transverse center lemma
  (Lemma~\ref{lem:transverse}) shows that this blow-up adds torsion to the
  fixed components only in a controlled way and creates no new fixed
  components (\S\ref{sec:ten}). This gives Theorem~\ref{thmA}: the first
  simply connected example below the stable Morse number in real dimension
  $<22$, and a counterexample to the Morse-number assertion of
  \cite{Ma2013} on a simply connected K\"ahler manifold of real dimension
  $10$.
\item In dimension $8$, necessary conditions together with the
  unconditional mechanism obstruction Theorem~\ref{thmB}, based on
  perfectness of moment maps over every prime field
  (Lemma~\ref{lem:perfect}, a specialization of the Borel--Smith
  inequality; we do not claim novelty for it).
\item A heuristic explanation of why the transverse-center construction
  needs real dimension at least $10$ (\S\ref{subsec:dimlimit}), and a
  machine classification of admissible torsion configurations
  corroborating Theorem~\ref{thmB} (Appendix~\ref{app:milp}).
\end{enumerate}
The low-dimensional statements of \S\ref{sec:six} are straightforward
consequences of Bai--Xu and duality. The K\"ahler examples in dimensions
$12$ and $16$ recorded in \S\ref{subsec:higher} are rearrangements of the
same ideas and are not claimed as separate results.

\subsection{Conventions}
Homology has integer coefficients unless indicated. For a finite abelian
group $A$ and a prime $p$, write $d_p(A)=\dim_{\F_p}(A\otimes\F_p)$; then
$d(A)=\max_p d_p(A)$. For a space $X$ set $t_i^p(X)=d_p(\Tor H_i(X))$ and
$r_i(X)=\rank H_i(X)$. Homology shifted by $2k$ is written $H(\cdot)[2k]$.

%======================================================================
\section{Torsion counts and fixed-point bounds}\label{sec:prelim}

\subsection{Degreewise and cyclic generator counts}

\begin{definition}\label{def:lambdaN}
For a closed manifold $W$ and $N\in\{1,2,\dots\}\cup\{\infty\}$ put
\[
  \lambda_N(W)=b(W)+2\sum_{c\in\Z/2N}
  d\Bigl(\bigoplus_{i\equiv c\ (2N)}\Tor H_i(W)\Bigr),
\]
with $\Z/2\infty=\Z$, so that $\lambda_\infty=\lambda$.
\end{definition}

Since $d(A\oplus B)=\max_p(d_p(A)+d_p(B))\le d(A)+d(B)$, one has
$\lambda_N\le\lambda$, with equality as soon as each residue class
mod $2N$ contains at most one degree carrying torsion. For finite groups
$A,B$,
\begin{equation}\label{eq:dsum}
  d(A)+d(B)-d(A\oplus B)>0\iff
  \operatorname{argmax}_p d_p(A)\cap\operatorname{argmax}_p d_p(B)=\varnothing
  \text{ and }A,B\ne0 .
\end{equation}
We say that $A$ and $B$ have \emph{different dominant primes} in this case.

\subsection{Morse-theoretic bounds}\label{subsec:morse}

\begin{proposition}[{\cite[Prop.~2.2]{OpenAI22}}]\label{prop:SMlower}
For every connected closed smooth manifold $W$, $\SM(W)\ge\lambda(W)$;
in particular every Morse function on $W$ has at least $\lambda(W)$
critical points.
\end{proposition}

\begin{theorem}[Smale]\label{thm:smale}
Let $W$ be a simply connected closed smooth manifold of dimension
$D\ge 6$. Then $W$ admits a Morse function with exactly
$r_i(W)+t_i(W)+t_{i-1}(W)$ critical points of index $i$; hence
$\SM(W)=\Morse(W)=\lambda(W)$.
\end{theorem}

For $D\ge 9$ this is \cite[Thm.~2.4]{OpenAI22} with the handle proof of
its appendix, based on \cite{Smale1963}. For all $D\ge6$ it is
\cite[Thm.~6.1]{Smale1962}.\footnote{The torsion-inclusive statement of
\cite[Thm.~6.1]{Smale1962} for $n\ge6$ was confirmed through secondary
sources (Encyclopedia of Mathematics, ``Morse inequalities'';
arXiv:2607.20259); the original wording was not consulted.}

\subsection{The cyclic integral Floer bound}

Let $N_M$ be the minimal Chern number of $(M,\omega)$: the positive
generator of $c_1(TM)(\pi_2(M))\subset\Z$, or $\infty$ if this subgroup
is zero (Bai--Xu write $N=0$ in this case).

\begin{theorem}[Bai--Xu {\cite[Thm.~A]{BaiXu2022}}, see also {\cite[Cor.~1.2]{BaiXu2026}}]\label{thm:BX}
Let $(M,\omega)$ be any closed symplectic manifold and $H$ a one-periodic
Hamiltonian with only nondegenerate fixed points. Then
\[
  \#\Fix_0(\phi_H^1;H)\ \ge\ \rank H_*(M;\Q)+2\sum_{i\in\Z/2N_M}\tau_i^{(2N_M)}(M)
  =\lambda_{N_M}(M),
\]
where $\tau_i^{(2N)}(M)$ is the number of nontrivial invariant factors of
the torsion of $\bigoplus_{k\equiv i\ (2N)}H_k(M;\Z)$, that is, its
minimal number of generators.
\end{theorem}

\begin{remark}\label{rem:rezchikov}
Rezchikov \cite[Thm.~1]{Rezchikov2022} phrases an integral lower bound as
the minimal rank of a cochain complex of free $\Z$-modules with cohomology
$H^*(M;\Z)$. His Floer complexes are built from cappings modulo the
symplectic area only, and he explicitly avoids the integrally graded
variant (\cite[\S2.1, Remark~2]{Rezchikov2022}: ``This allows one to define
integrally graded Floer groups, which we will avoid in this paper''), so the
grading is cyclic and the bound obtained is the cyclic one,
consistent with Theorem~\ref{thm:BX}. Read with a $\Z$-grading, the
statement would give $\#\Fix\ge\lambda(M)$, which Theorem~\ref{thmA}
contradicts.
\end{remark}

\subsection{Perturbation of finite-order Hamiltonian symmetries}

The following is \cite[Prop.~4.2]{OpenAI22} (Prop.~dyn:involution there)
with the involution replaced by an element of finite order.

\begin{proposition}\label{prop:perturb}
Let $(P,\omega)$ be a closed symplectic manifold and $g$ a Hamiltonian
diffeomorphism with $g^k=\id$ for some $k\ge1$. Let
$\Fix(g)=C_1\sqcup\dots\sqcup C_r$ and let $h_j\colon C_j\to\R$ be Morse
functions. There is a smooth $g$-invariant $f\colon P\to\R$ with
$f|_{C_j}=h_j$ such that for all sufficiently small $\epsilon>0$ the
Hamiltonian flow $\psi^s$ of $f$ satisfies
\[
  \Fix(\psi^\epsilon\circ g)=\bigsqcup_j\Crit(h_j),
\]
every such fixed point is nondegenerate, and $\psi^\epsilon\circ g$ is the
time-one map of a smooth one-periodic Hamiltonian. If $P$ is simply
connected, all fixed-point orbits are contractible.
\end{proposition}

\begin{proof}
We indicate the changes to the proof of \cite[Prop.~4.2]{OpenAI22}.
Average a metric over $\langle g\rangle$; the fixed set is a union of
closed submanifolds with $T_xC_j=E_x^+=\ker(d_xg-I)$. Let $E_x^{\ne}$ be
the sum of the (real parts of the) generalized eigenspaces of $d_xg$ for
eigenvalues $\ne1$; $d_xg$ is semisimple of finite order, so
$T_xP=E_x^+\oplus E_x^{\ne}$ and $E_x^{\ne}=\operatorname{im}(d_xg-I)$.
If $u\in E_x^+$ and $v=d_xg(w)-w$, then
$\omega(u,v)=\omega(d_xg\,u,d_xg\,w)-\omega(u,w)=0$, so the summands are
$\omega$-orthogonal and $C_j$ is symplectic. Let
$f=\frac1k\sum_{m=0}^{k-1}\widetilde f\circ g^m$ for an extension
$\widetilde f$ of the $h_j$. Then $f$ is $g$-invariant, $d_xf$ vanishes on
$\operatorname{im}(d_xg-I)=E_x^{\ne}$ for $x\in\Fix(g)$, hence $X_f$ is
tangent to $C_j$ and equals $X_{h_j}$ there, and $g$ commutes with the flow
$\psi^s$ of $X_f$. Put $\Phi=\psi^\epsilon\circ g$; then
$\Phi^k=\psi^{k\epsilon}$. If $\Phi(x)=x$ and $k\epsilon$ is smaller than
the constant of \cite[Lemma~4.1]{OpenAI22} (short-period lemma), then
$X_f(x)=0$, so $\psi^\epsilon(x)=x$ and $g(x)=x$, and $x\in\Crit(h_j)$.
At such $x$, with $A=d_xX_f$ preserving $E^+_x$ and $E^{\ne}_x$, the tangent
block of $d_x\Phi-I$ is $\exp(\epsilon A^+)-I=\epsilon A^++O(\epsilon^2)$
with $A^+$ invertible by the Morse condition, and the normal block is
$d_xg|_{E^{\ne}}\exp(\epsilon A^{\ne})-I\to d_xg|_{E^{\ne}}-I$, which is
invertible. The concatenation argument producing a one-periodic
Hamiltonian is unchanged.
\end{proof}

\subsection{Integral blow-up formula}

\begin{proposition}[{\cite[Thm.~4.1]{DuanLi2009}}, {\cite[Prop.~3.3]{OpenAI22}}]\label{prop:blowup}
Let $A$ be a compact connected K\"ahler manifold and $Z\subset A$ a closed
connected complex submanifold of complex codimension $c\ge2$. Then
$\Bl_ZA$ is compact K\"ahler, simply connected if $A$ is, and
\[
  H_i(\Bl_ZA;\Z)\cong H_i(A;\Z)\oplus\bigoplus_{j=1}^{c-1}H_{i-2j}(Z;\Z).
\]
\end{proposition}

%======================================================================
\section{Low dimensions}\label{sec:six}

The statements of this section are direct consequences of
Theorem~\ref{thm:BX} and Poincar\'e duality; we record them to delimit
Corollary~\ref{cor:mindim}.

\emph{Real dimension $2$.} For a closed oriented surface $\Sigma_g$,
$\SM=\Morse=\lambda=2g+2$ (Proposition~\ref{prop:SMlower} and the standard
Morse function). If $g=0$, $N=2$ and $\lambda_2(S^2)=2$; if $g\ge1$,
$\pi_2=0$, $N=\infty$ and $\lambda_\infty=\lambda$. Theorem~\ref{thm:BX}
gives $\#\Fix_0\ge2g+2=\SM(\Sigma_g)$, so there is no counterexample.

\emph{Simply connected real dimension $4$.} The homology is torsion-free,
so Theorem~\ref{thm:BX} gives $\#\Fix\ge b(M)=\lambda(M)$. A simply
connected four-dimensional counterexample would therefore require
$\SM(M)>\lambda(M)$. Smale's theorem does not apply in dimension $4$, and
we do not settle whether $\SM=\lambda$ for simply connected closed
$4$-manifolds here.

\begin{proposition}\label{prop:six}
Let $(M,\omega)$ be a simply connected closed symplectic manifold of real
dimension $6$. Every one-periodic Hamiltonian with only nondegenerate fixed
points satisfies $\#\Fix(\phi_H^1)\ge\lambda(M)=\SM(M)$.
\end{proposition}

\begin{proof}
Write $T_i=\Tor H_i(M)$. By universal coefficients and Poincar\'e duality,
$T_i\cong\Tor H^{i+1}(M)\cong T_{5-i}$. Since $H_0$ is free and $H_1=0$,
we get $T_0=T_1=T_4=T_5=0$, $T_6=0$, and $T_2\cong T_3$. The two degrees
$2$ and $3$ have different parities, hence lie in different residue
classes modulo $2N$ for every $N\in\{1,2,\dots,\infty\}$. Therefore
$\lambda_N(M)=\lambda(M)$ for every $N$. Theorem~\ref{thm:BX} gives
$\#\Fix\ge\lambda(M)$, and Theorem~\ref{thm:smale} with $D=6$ gives
$\lambda(M)=\SM(M)$; contractibility is automatic.
\end{proof}

\begin{remark}
Proposition~\ref{prop:six} concerns nondegenerate Hamiltonians. The
six-dimensional example of \cite{OpenAIDegenerate}, a Hamiltonian
diffeomorphism of the quadric $Q^3$ with three fixed points, has a
degenerate fixed point and concerns the critical-point count; it does not
conflict with Proposition~\ref{prop:six}.
\end{remark}

%======================================================================
\section{An Artin--Mumford quotient threefold}\label{sec:F}

\subsection{An invariant quartic symmetroid web}\label{subsec:web}
Let $\zeta=e^{2\pi i/5}$ and let $G=\langle g\rangle\cong\Z/5$ act on
$\PP^3$ by
\[
  g\cdot[x_0:x_1:x_2:x_3]=[x_0:\zeta^4x_1:\zeta^3x_2:\zeta x_3].
\]
Consider the web of quadrics given by the symmetric matrix of linear forms
\[
  A(x)=\begin{pmatrix}
  6x_0&5x_1&4x_2&0\\
  5x_1&8x_2&0&4x_3\\
  4x_2&0&7x_3&6x_0\\
  0&4x_3&6x_0&5x_1
  \end{pmatrix},
\]
and let $B=\{f=0\}$, $f=\det A$:
\[
  f=900x_0^2x_1^2-1728x_0^3x_2-640x_1x_2^3-875x_1^3x_3+720x_0x_1x_2x_3
  +256x_2^2x_3^2-672x_0x_3^3 .
\]
Every monomial of $f$ has $g$-weight $3\bmod5$, so $f(gx)=\zeta^3f(x)$.

\begin{lemma}[machine-verified, Appendix~\ref{app:sage}]\label{lem:webfacts}
\begin{enumerate}
\item The four matrices $A_0,\dots,A_3$ (coefficients of $x_0,\dots,x_3$)
  are linearly independent, the web has no base points, and no member has
  rank $\le1$.
\item $f$ is irreducible over $\Q$, the singular scheme of $B$ equals the
  rank-$\le2$ scheme of $A$, and it is reduced of length $10$. Thus $B$ has
  exactly ten singular points, all ordinary double points ($A_1$). Since
  $\operatorname{Sing}(B)$ is finite and a reducible surface is singular
  along a curve, $f$ is absolutely irreducible. The only $g$-fixed points
  of $\PP^3$ are the coordinate points, which are smooth points of $B$
  (item 3) and $|G|=5$ is prime, so the nodes form two free $G$-orbits.
\item The four coordinate points $p_0,\dots,p_3$ lie on $B$ and are smooth
  points of $B$.
\item (Beauville's condition (ii).) If a line of $\PP^3$ is contained in
  the singular locus of a quadric of the web, it is not contained in any
  other quadric of the web. Equivalently, at each of the ten rank-$2$
  points the three quadratic conditions on the web imposed by the kernel
  line are independent; this is checked exactly over $\overline{\Q}$.
\item The ten nodes impose independent conditions on quadrics: the
  Hilbert function of $\C[x_0,\dots,x_3]/I_{\mathrm{nodes}}$ is
  $1,4,10,10,\dots$, so $H^0(\PP^3,I_{\mathrm{nodes}}(2))=0$.
\end{enumerate}
\end{lemma}

Let $\overline X\to\PP^3$ be the double cover $w^2=f$ (a hypersurface in
$\PP(1,1,1,1,2)$) and $X:=\mathrm{AM}_5$ its blow-up at the ten nodes. The
lift $\hat g\colon(x,w)\mapsto(gx,\zeta^4w)$ has order $5$ (note
$\zeta^8=\zeta^3$) and lifts to $X$; we denote it again by $g$.

\begin{lemma}\label{lem:fixedpts}
$g$ acts on $X$ with exactly four fixed points, one over each $p_k$, each
of type $\frac15(1,1,4)$. All other points have trivial stabilizer.
\end{lemma}

\begin{proof}
The weights $0,4,3,1$ are pairwise distinct, so $\Fix(g|_{\PP^3})=
\{p_0,\dots,p_3\}$; the nodes form free orbits, so no exceptional quadric
contains a fixed point. Since $p_k\in B$, the fiber over $p_k$ is the single
point $w=0$. At $p_0$ use affine coordinates $u_k=x_k/x_0$ of weights
$4,3,1$. The linear part of $f(1,u)$ is a nonzero multiple of $u_2$
($\partial f/\partial x_2(p_0)=-1728$); hence $(w,u_1,u_3)$ are local
coordinates on $X$, of weights $(4,4,1)$, i.e.\ type $\frac15(1,1,4)$
after multiplying by $4$. At $p_1,p_2,p_3$ the nonvanishing partials are
$\partial_{x_3}f(p_1)=-875$, $\partial_{x_1}f(p_2)=-640$,
$\partial_{x_0}f(p_3)=-672$, and the same computation gives weights
$(1,1,4)$, $(3,2,3)$, $(2,3,2)$, all of type $\frac15(1,1,4)$. As $|G|$ is
prime, every non-fixed point has trivial stabilizer.
\end{proof}

\subsection{Topology of the invariant double solid}

\begin{proposition}\label{prop:AM5}
$X=\mathrm{AM}_5$ is simply connected and
\[
  H_*(X;\Z)=(\Z,\ 0,\ \Z^{11}\oplus\Z/2,\ \Z/2,\ \Z^{11},\ 0,\ \Z)
  \quad\text{in degrees }0,\dots,6.
\]
\end{proposition}

\begin{proof}
\emph{Betti numbers.} By Lemma~\ref{lem:webfacts}(5) the ten nodes impose
independent conditions on quadrics. For a double solid branched along a
surface of degree $d$ whose only singularities are $\mu$ nodes, Clemens'
defect $\delta$ is the number of dependent conditions imposed by the nodes
on forms of degree $\tfrac32d-4$, and the big resolution $\widetilde Y$
(blow-up of the nodes) has
\[
  h^{1,1}(\widetilde Y)=1+\mu+\delta,\qquad
  h^{1,2}(\widetilde Y)=\tbinom{3d/2-1}{3}-4\tbinom{d/2}{3}-\mu+\delta
\]
\cite{Clemens1983}; we use the formulation of \cite[\S1 and
Cor.~5.1]{Cynk2008}, which states these formulas verbatim (we have not
consulted the page of \cite{Clemens1983} where they appear). Here
$d=4$, $\mu=10$, and the forms have degree $2$, so $\delta=0$ and
$h^{1,1}(X)=11$, $h^{1,2}(X)=10-0-10+0=0$. Since $h^{3,0}=h^{2,0}=0$
for a quartic double solid, $b_3(X)=0$ and $b_2(X)=11$. This agrees with
the Euler characteristic: a quartic with ten nodes has $\chi=24-10=14$,
so $\chi(\overline X)=2\chi(\PP^3)-14=-6$, and replacing each node by a
quadric adds $3$, giving $\chi(X)=24=2+2\cdot11$.

\emph{Nonzero $2$-torsion and simple connectivity.} The web is base point
free and satisfies Beauville's condition (ii)
(Lemma~\ref{lem:webfacts}(1),(4)). These are exactly the hypotheses of
\cite[\S6.3]{Beauville2015}. Under them Beauville shows that the
discriminant has ten ordinary double points and no other singularities,
and \cite[Thm.~6 and its proof]{Beauville2015} shows that the blown-up
double solid $X$ is unirational and that $\Tor H^3(X;\Z)$ contains a
nonzero element of order $2$. Unirational smooth projective varieties are
simply connected (Serre), so $\pi_1(X)=1$.

\emph{Exactness, by deformation.} Beauville's argument gives
$\Tor H^3\ne0$, and Voisin \cite[proof of Cor.~2.6]{Voisin2015}
observes that every nodal quartic double solid has a unirational
parametrization of degree $2$, so the torsion is killed by $2$. We could
not find a source stating that $\Tor H^3\cong\Z/2$ (rather than
$(\Z/2)^k$) for \emph{every} web satisfying (i) and (ii). Endra{\ss}
\cite[Lemma~1.2]{Endrass1999} embeds the torsion into a group of even
sets of nodes, but we have not computed that group for our web.
We therefore deform to the Artin--Mumford example, for which
$\Tor H^3\cong\Z/2$ \cite{ArtinMumford1972}, \cite[Constr.~2.2]{Deligne1973}.
Since $\Tor H^3$ is a birational invariant of smooth projective
threefolds \cite[Prop.~2.1]{Deligne1973}, this holds for every smooth
projective model of that example, in particular for the blown-up double
solid of the corresponding symmetroid (\cite[\S6.3]{Beauville2015}).
\todo{confirm that the Artin--Mumford example of \cite{Deligne1973} is
birational to a double solid branched along a quartic symmetroid
satisfying (i) and (ii); \cite[\S6.3]{Beauville2015} and Cossec state that
the Artin--Mumford example is such a double solid.}

Let $\Gr=\Gr(4,\Sym^2\C^4)$ (irreducible) and let
$\mathcal D\subset\PP^3\times\Gr$ be the relative rank-$\le2$ scheme and
$\mathcal S\supset\mathcal D$ the relative singular scheme of the
symmetroids; $\mathcal S$ is proper over $\Gr$. The partial derivatives of
$\det A_\Lambda$ are combinations of $3\times3$ minors, so
$\mathcal D_\Lambda\subset\mathcal S_\Lambda$. The variety $\Sigma_2\subset
\PP^9$ of symmetric $4\times4$ matrices of rank $\le2$ is irreducible,
arithmetically Cohen--Macaulay, of codimension $3$ and degree $10$; the
degree is the case $n=4$, $r=2$ of the Harris--Tu formula
$\prod_{j=0}^{n-r-1}\binom{n+j}{n-r-j}/\binom{2j+1}{j}$
\cite{HarrisTu1984} (we use the statement as quoted in
\cite[\S1]{BHORS2012}). When $\mathcal D_\Lambda$ is finite, it is a proper
linear section of the Cohen--Macaulay variety $\Sigma_2$, so its length
equals $\deg\Sigma_2=10$. Hence $\operatorname{length}\mathcal S_\Lambda\ge10$
whenever it is finite. On the open locus over which $\mathcal S$ is
quasi-finite, the fibre length is upper semicontinuous, and the \'etale
locus of a finite flat family is open; hence
\[
  U_0=\{\Lambda:\mathcal S_\Lambda\text{ is reduced of length }10\}
\]
is Zariski open in $\Gr$. It contains every web satisfying (i) and (ii)
(by Beauville's statement above), in particular the Artin--Mumford web and
our web (Lemma~\ref{lem:webfacts}(2)). Being a nonempty Zariski open subset
of an irreducible variety, it is connected. Over $U_0$ the nodes form a
finite \'etale multisection, and by the holomorphic Morse lemma with
parameters the double solids are locally analytically of the form
$w^2=y_1^2+y_2^2+y_3^2$ times the base near the nodes. Blowing up the
multisection therefore gives a smooth proper family, and by Ehresmann's
theorem all members are diffeomorphic. Therefore
$\Tor H^3(X;\Z)\cong\Z/2$. Universal coefficients and duality give
$\Tor H_2\cong\Tor H^3=\Z/2$ and $H_3\cong H^3=\Z/2$.
\end{proof}

\subsection{Quotient, resolution, fundamental group}

Let $X^\circ\subset X$ be the complement of small $G$-invariant open balls
around the four fixed points, $U=X^\circ/G$, and $\partial U=\bigsqcup_{k=0}^3L_k$
with $L_k\cong S^5/(\Z/5)$ lens spaces. Each fixed point has type
$\frac15(1,1,4)$; let $\rho_k\colon R_k\to\C^3/G$ be the toric resolution
given by iterated star subdivision of the cone of $\frac15(1,1,4)$ at
the lattice points $\frac15(1,1,4),\frac15(2,2,3),\frac15(3,3,2),
\frac15(4,4,1)$ (Appendix~\ref{app:sage}). It has four exceptional
divisors and nine maximal cones; it is a composition of weighted blow-ups,
hence projective. Define
\[
  F=U\cup_{\partial U}\textstyle\bigsqcup_kR_k
  =\text{the resolution of }X/G\text{ at its four singular points.}
\]
$F$ is a smooth projective threefold. Each $R_k$, a small neighbourhood
of the exceptional set of $\rho_k$, deformation retracts onto the
exceptional set $E=D_1\cup D_2\cup D_3\cup D_4$, where $D_i$ is the toric
divisor of the interior ray $v_i=\frac i5(1,1,4)\bmod\Z^3$. In the fan
(computed in \texttt{computations/sage/bF2.sage}) the maximal cones are
$\langle e_1,e_2,v_4\rangle$, $\langle e_3,e_1,v_1\rangle$,
$\langle e_3,e_2,v_1\rangle$ and $\langle e_a,v_i,v_{i+1}\rangle$
($a=1,2$, $i=1,2,3$). Hence each $D_i$ is a smooth complete toric surface,
$D_i\cap D_j=\varnothing$ for $|i-j|\ge2$, there are no triple
intersections, and $D_i\cap D_{i+1}$ is the invariant curve
$V(\langle v_i,v_{i+1}\rangle)\cong\CP^1$. A smooth complete toric surface
has free homology concentrated in even degrees \cite[\S5.2,
pp.~101--103]{Fulton1993}. An invariant curve on it is primitive in $H_2$,
since it meets an adjacent invariant curve with intersection number $1$.
Mayer--Vietoris along the chain $D_1\cup D_2\cup D_3\cup D_4$ therefore
shows that $H_*(R_k)\cong H_*(E)$ is free and concentrated in even
degrees. Its Betti numbers are $(1,4,4)$ (the Euler characteristic $9$ is
the number of maximal cones).

\begin{lemma}\label{lem:pi1F}
$\pi_1(F)=1$.
\end{lemma}

\begin{proof}
$X^\circ$ is simply connected (removing points from a simply connected
$6$-manifold), so $\pi_1(U)=G$ since $G$ acts freely on $X^\circ$. A loop in
$L_k$ lifting to a path from $x$ to $gx$ in the boundary sphere represents
$g$, so $\pi_1(L_k)\to\pi_1(U)$ is onto. The $R_k$ are simply connected.
Van Kampen gives $\pi_1(F)=1$. (Equivalently, by Armstrong's theorem
\cite{Armstrong1968} $\pi_1(X/G)=G/\langle\text{elements with fixed
points}\rangle=1$, and resolving isolated quotient singularities with
simply connected exceptional sets does not change $\pi_1$.)
\end{proof}

\subsection{Integral homology of $F$}\label{subsec:Fhom}

\begin{proposition}\label{prop:Fhom}
$H_1(F)=H_5(F)=0$, $b_1=b_3=b_5=0$, $b_2(F)=b_4(F)=19$, $b(F)=40$, and
$\Tor H_i(F)\cong\Tor H_{5-i}(F)$; in particular $\Tor H_4(F)=0$.
\end{proposition}

\begin{proof}
By the Lefschetz formula
$\chi(X/G)=\frac15\bigl(\chi(X)+4\cdot\#\Fix(g)\bigr)=\frac15(24+16)=8$.
Each resolution replaces a point by a set of Euler characteristic $9$,
so $\chi(F)=8+4\cdot8=40$. Rationally
$H_*(U;\Q)=H_*(X^\circ;\Q)^G$ and $H_{\mathrm{odd}}(X;\Q)=0$; with the
even homology of the $R_k$, Mayer--Vietoris gives $b_{\mathrm{odd}}(F)=0$,
hence $b_2=b_4=19$ (the invariant part of $H_2(X;\Q)$ has rank $3$ by the
permutation-module description in Lemma~\ref{lem:torF}, and
$3+4\cdot4=19$). $H_1(F)=0$ by Lemma~\ref{lem:pi1F}. $F$ is a closed
oriented $6$-manifold, so $\Tor H_i\cong\Tor H^{i+1}\cong\Tor H_{5-i}$.
\end{proof}

\begin{lemma}\label{lem:torF}
$\Tor H_2(F)\cong\Tor H_3(F)\cong\Z/2$. In particular $F$ has no
$5$-torsion.
\end{lemma}

\begin{proof}
By Proposition~\ref{prop:Fhom} it suffices to compute $\Tor H_2(F)$ and
$\Tor H_3(F)$ one prime at a time. Mayer--Vietoris for $F=U\cup R$,
$R=\bigsqcup R_k$, reads
\begin{equation}\label{eq:MV}
  H_i(\partial U)\to H_i(U)\oplus H_i(R)\to H_i(F)\to H_{i-1}(\partial U),
\end{equation}
with $H_*(L_k)=(\Z,\Z/5,0,\Z/5,0,\Z)$ and $H_*(R)$ free, in even degrees.

\emph{Step 1 (input).} $\Tor H_2(X)=\Z/2$ and $H_3(X)=\Z/2$ exactly
(Proposition~\ref{prop:AM5}), and $H_i(X^\circ)=H_i(X)$ for $i\le4$.

\emph{Step 2 (localization at $\ell\ne5$).} $H_i(\partial U)_{(\ell)}=0$
for $i=1,2,3$, so \eqref{eq:MV} gives
$H_i(F)_{(\ell)}\cong H_i(U)_{(\ell)}\oplus H_i(R)_{(\ell)}$ for $i=2,3$,
and $\Tor H_i(F)_{(\ell)}=\Tor H_i(U)_{(\ell)}$.

\emph{Step 3 (transfer).} As $|G|$ is invertible in $\Z_{(\ell)}$, the
Cartan--Leray spectral sequence collapses to
$H_i(U;\Z_{(\ell)})\cong H_i(X^\circ;\Z_{(\ell)})_G\cong
H_i(X^\circ;\Z_{(\ell)})^G$ (averaging).

\emph{Step 4 (torsion of invariants).} For a $\Z_{(\ell)}[G]$-module $H$
with $|G|$ invertible, $H^1(G;\Tor H)=0$, so $\Tor(H^G)=(\Tor H)^G$. For
$\ell=2$, $G$ acts trivially on $\Z/2$, so $\Tor H_2(F)_{(2)}\cong
\Tor H_3(F)_{(2)}\cong\Z/2$. For $\ell\notin\{2,5\}$ there is no
$\ell$-torsion.

\emph{Step 5 (the prime $5$).} Let $h\in H^2(X)$ be the pullback of the
hyperplane class and $E_1,\dots,E_{10}$ the classes of the exceptional
quadrics; $g$ fixes $h$ and permutes the $E_i$ in two free orbits. Pair
them with the curve classes $\tilde\ell$ (preimage of a general line,
$h\cdot\tilde\ell=2$, $E_i\cdot\tilde\ell=0$) and $f_i$ (a ruling of the
$i$-th quadric, $h\cdot f_i=0$, $E_j\cdot f_i=-\delta_{ij}$, since
$E_i|_{E_i}=\mathcal O(-1,-1)$). The intersection matrix is
$\operatorname{diag}(2,-1,\dots,-1)$, of determinant $\pm2$. Hence
$h,E_1,\dots,E_{10}$ span a sublattice of index dividing $2$ in
$H^2(X)/\Tor$, and
\[
  H^2(X;\Z_{(5)})\cong H_2(X;\Z_{(5)})\cong\Z_{(5)}\oplus\Z_{(5)}[G]^2
\]
as $\Z_{(5)}[G]$-modules (permutation modules are self-dual). In the
Cartan--Leray spectral sequence $E^2_{p,q}=H_p(G;H_q(X^\circ;\Z_{(5)}))
\Rightarrow H_{p+q}(U;\Z_{(5)})$ \cite[VII.7]{Brown1982} we have
$H_1(X^\circ)=0$, $H_3(X^\circ;\Z_{(5)})=0$, and
$E^2_{0,2}=(\Z\oplus\Z[G]^2)_G\otimes\Z_{(5)}\cong\Z_{(5)}^3$. Since
$H_2(G;\Z)=0$, the only possible differential into $E_{0,2}$ comes from
the finite group $E^3_{3,0}=H_3(G;\Z)=\Z/5$; it is zero because
$E^2_{0,2}\cong\Z_{(5)}^3$ is torsion free. Hence
$H_2(U;\Z_{(5)})\cong\Z_{(5)}^3$. Since moreover $b_3(U)=0$, universal
coefficients give $H^3(U;\Z)\otimes\Z_{(5)}=0$, and Lefschetz duality
($U$ is a compact oriented $6$-manifold with boundary) gives
$H_3(U,\partial U)_{(5)}\cong H^3(U)_{(5)}=0$. Hence
$H_3(\partial U)\to H_3(U)$ is onto after localizing at $5$. As
$H_3(R)=0$ and $H_2(\partial U)=0$, \eqref{eq:MV} gives
$H_3(F)_{(5)}=\operatorname{coker}(H_3(\partial U)\to H_3(U))_{(5)}=0$.

\emph{Step 6 (duality at $5$).} From \eqref{eq:MV} alone, $H_2(F)_{(5)}$
is an extension of $\ker\bigl(H_1(\partial U)\to H_1(U)\bigr)\cong(\Z/5)^3$
by the free module $H_2(U)_{(5)}\oplus H_2(R)_{(5)}$, so Mayer--Vietoris
does not by itself exclude a $\Z/5$ in $H_2(F)$. Torsion-freeness comes
from duality: $\Tor H_2(F)_{(5)}\cong\Tor H_3(F)_{(5)}=0$ by Step~5.
\end{proof}

\begin{remark}\label{rem:5torsion}
The exactness of Lemma~\ref{lem:torF} is essential. If
$\Tor H_2(F)=(\Z/2)^2$, or if $F$ had $5$-torsion in degrees $2,3$, the
gap of Theorem~\ref{thmA} (and of Remark~\ref{rem:dim12}) would collapse
to $\le0$ (\texttt{computations/homology/w10.py},
\texttt{computations/homology/lam.py}). Note that $H_3(U;\Z_{(5)})$ itself
has order $25$ (contributions $E^\infty_{1,2}$ and $E^\infty_{3,0}$); it is
killed in $F$ by the lens space boundary.
\end{remark}

\subsection{A Godeaux-type divisor}\label{subsec:S}

The $G$-invariant quintic monomials (weight sum $\equiv0\bmod5$) span a
$12$-dimensional space. Fix the invariant quintic
\[
\begin{aligned}
  Q={}&6x_0^5+9x_1^5+17x_0x_1^3x_2+15x_0^2x_1x_2^2+13x_2^5+14x_0^3x_1x_3
  +10x_1^2x_2^2x_3\\
  &+10x_0x_2^3x_3+x_0x_1^2x_3^2+19x_0^2x_2x_3^2+15x_1x_2x_3^3+5x_3^5 .
\end{aligned}
\]

\begin{lemma}[machine-verified, Appendix~\ref{app:sage}]\label{lem:quintic}
The surface $\{Q=0\}$ is smooth, avoids $p_0,\dots,p_3$ and the ten nodes
of $B$, and meets $B$ transversally in a smooth curve $\Delta$.
\end{lemma}

Let $\widetilde S\subset X$ be the preimage of $\{Q=0\}$; it is the double
cover of the quintic surface branched along $\Delta\in|4H|$, and lies in the
locus where $X\to\overline X$ is an isomorphism and $G$ acts freely. Put
$S=\widetilde S/G\subset F$.

\begin{proposition}\label{prop:S}
$S$ is a smooth divisor in $F$ with $\pi_1(S)=\Z/5$; its homology is
\[
  H_*(S;\Z)=(\Z,\ \Z/5,\ \Z^{40}\oplus\Z/5,\ 0,\ \Z).
\]
\end{proposition}

\begin{proof}
Smoothness and the divisor property follow from Lemma~\ref{lem:quintic}
and the fact that $\widetilde S$ avoids the fixed points. Simple
connectivity of $\widetilde S$: by Nori's theorem \cite{Nori1983}, since
$\Delta$ is a smooth curve with $\Delta^2>0$ on the simply connected
quintic surface, $\pi_1$ of the complement of $\Delta$ is abelian, hence
cyclic generated by a meridian $\gamma$; the double cover corresponds to
$\langle\gamma^2\rangle$, which is generated by a meridian of
$\widetilde\Delta$ and dies in $\widetilde S$. Here
$H_1(\Sigma\setminus\Delta)=\Z/4$, because the hyperplane class $H$ of the
quintic surface $\Sigma$ is primitive ($H^2=5$ is squarefree), so
$\Delta\in|4H|$ has divisibility exactly $4$. The abelianness statement is
\cite[Prop.~3.27]{Nori1983}, applied with $D=\Delta$ (smooth, so
$r(\Delta)=0$ and the hypothesis $\Delta^2>2r(\Delta)$ holds) and
$E=\varnothing$: the kernel of $\pi_1(\Sigma\setminus\Delta)\to\pi_1(\Sigma)
=1$ is finitely generated abelian.
Hence $\pi_1(S)=G=\Z/5$, $H_1(S)=\Z/5$, $\Tor H_2(S)\cong\Tor H_1(S)=\Z/5$
and $H_3(S)=0$. For Betti numbers, $\chi$ of a smooth quintic surface is
$55$ and $\Delta$ has $2g-2=\Delta\cdot(\Delta+K)=4H\cdot5H=100$, so
$\chi(\widetilde S)=2\cdot55+100=210$, $\chi(S)=42$ and $b_2(S)=40$.
\end{proof}

We call $S$ \emph{Godeaux-type} because $\pi_1(S)=\Z/5$; it is not a
Godeaux surface (its geometric genus is positive).

\begin{remark}
This gives, in the K\"ahler category, a simply connected compact K\"ahler threefold with $2$-primary
Brauer-type torsion containing a smooth divisor with $5$-torsion in $H_1$.
The Brauer group is a birational invariant, which is why one cannot start
from a rational threefold such as the toric resolution of $\PP^3/\Z_5$.
\end{remark}

%======================================================================
\section{The ten-dimensional construction}\label{sec:ten}

\subsection{The surface $Y$ and two invariant lines}

Let $Y$ be $\CP^1\times\CP^1$ blown up at the four points
$\{0,\infty\}^2$, with the diagonal circle action
$t\cdot(z,w)=(e^{2\pi it}z,e^{2\pi it}w)$ lifted to $Y$ and an invariant
K\"ahler form. By \cite[Prop.~4.3]{OpenAI22}, the half-turn $g_Y$ is a
Hamiltonian involution with $\Fix(g_Y)=E_{00}\sqcup E_{0\infty}\sqcup
E_{\infty0}\sqcup E_{\infty\infty}$ (the exceptional curves), $dg_Y=+1$ on
$TE$ and $-1$ on the normal line; $Y$ is simply connected with
$H_*(Y)=(\Z,0,\Z^6,0,\Z)$; and $c_1(TY)[E]=1$.

Let $L_0,L_2\subset Y$ be the proper transforms of $\{0\}\times\CP^1$ and
$\{\infty\}\times\CP^1$. They are disjoint invariant smooth rational
curves; $g_Y$ rotates each by $\pi$, with fixed points (``poles'')
$L_0\cap E_{00}$, $L_0\cap E_{0\infty}$, $L_2\cap E_{\infty0}$,
$L_2\cap E_{\infty\infty}$. Thus every fixed curve contains exactly one
pole, and at a pole $y$ the curves $L$ and $E$ meet transversally, so
$T_yY=T_yE\oplus T_yL$ with $dg_Y=\operatorname{diag}(1,-1)$.

\subsection{Transverse center lemma}

\begin{lemma}[transverse center lemma]\label{lem:transverse}
Let $A$ be a complex manifold, $g$ a holomorphic automorphism of finite
order, $\Fix(g)=\bigsqcup_jC_j$, and $W\subset A$ a $g$-invariant closed
complex submanifold of codimension $\ge2$. Assume that for every
$p\in W\cap\Fix(g)$
\begin{equation}\label{eq:transverse}
  E_p^{\ne}:=\operatorname{im}(d_pg-I)\subset T_pW .
\end{equation}
Let $\hat g$ be the lift of $g$ to $\widehat A=\Bl_WA$. Then
\[
  \Fix(\hat g)=\bigsqcup_j\widehat C_j,\qquad
  \widehat C_j\cong\Bl_{W\cap C_j}C_j ,
\]
where $\widehat C_j$ is the proper transform of $C_j$ (equal to $C_j$ if
$W\cap C_j=\varnothing$, and to $C_j$ blown up along a smooth center
otherwise). In particular $\hat g$ has no fixed components other than the
proper transforms of those of $g$.
\end{lemma}

\begin{proof}
Write $\pi\colon\widehat A\to A$ and $E=\pi^{-1}(W)=\PP(N_W)$.

\emph{(a) Cleanness.} $W\cap\Fix(g)=\Fix(g|_W)$ is a closed submanifold
of $W$ with $T_p(W\cap C_j)=\ker(d_pg-I)\cap T_pW=T_pC_j\cap T_pW$; so $W$
and each $C_j$ intersect cleanly. By \eqref{eq:transverse},
$T_pA=T_pC_j\oplus E_p^{\ne}\subset T_pC_j+T_pW$, so
$T_pA=T_pC_j+T_pW$.

\emph{(b) Fixed points off $E$, and over non-fixed points of $W$.} $\pi$ is a
$g$-equivariant isomorphism off $E$, so $\Fix(\hat g)\setminus E=
\Fix(g)\setminus W$. If $p\in W$ and $g(p)\ne p$, then $\hat g$ maps the
fiber $E_p$ to $E_{g(p)}\ne E_p$, so it has no fixed points there.

\emph{(c) Fibers over $W\cap\Fix(g)$.} For $p\in W\cap C_j$, $\hat g$ acts on
$E_p=\PP(T_pA/T_pW)$ through the map induced by $d_pg$. By
\eqref{eq:transverse}, $T_pA/T_pW$ is the image of $E_p^+=T_pC_j$, on
which $d_pg=\id$. Hence $\hat g$ fixes $E_p$ pointwise.

\emph{(d) Identification.} By (a) the natural map
$T_pC_j/T_p(W\cap C_j)\to T_pA/T_pW$ is injective (cleanness) and surjective
($T_pA=T_pC_j+T_pW$), so $\PP(N_{W\cap C_j/C_j})|_p\to E_p$ is an
isomorphism. By cleanness the scheme-theoretic intersection $W\cap C_j$ is
the smooth submanifold above, so the strict transform $\widehat C_j$
(the closure of $C_j\setminus W$) is isomorphic to $\Bl_{W\cap C_j}C_j$
\cite[II.7.15]{Hartshorne1977}; it is a closed connected submanifold whose
intersection with $E$ is its exceptional divisor
$\PP(N_{W\cap C_j/C_j})$, which by the isomorphism just noted is all of
$E|_{W\cap C_j}$. It is contained in the closed set $\Fix(\hat g)$.
Combining (b) and (c), as sets
\[
  \Fix(\hat g)=\bigl(\Fix(g)\setminus W\bigr)\cup E|_{W\cap\Fix(g)}
  =\bigsqcup_j\bigl((C_j\setminus W)\cup E|_{W\cap C_j}\bigr)
  =\bigsqcup_j\widehat C_j .
\]
The $\widehat C_j$ are disjoint because the $C_j$ are, and $\Fix(\hat g)$,
a closed complex submanifold, is their disjoint union.
\end{proof}

\begin{remark}
Without \eqref{eq:transverse} the lemma fails: blowing up $S\times\{y\}$
with $y$ a pole, or $S\times E$, produces additional fixed components
(copies of $S$ or $S\times\CP^1$ inside $E$) that carry torsion and destroy
the gap; this is a direct computation with the bookkeeping of
\texttt{computations/homology/lam.py}. See also
\S\ref{subsec:dimlimit}.
\end{remark}

\subsection{The manifold $M$ and its fixed locus}

Let $F$ and $S\subset F$ be as in \S\ref{sec:F} and put
\[
  W=S\times(L_0\sqcup L_2)\subset F\times Y,\qquad
  M=\Bl_W(F\times Y).
\]
$W$ has complex dimension $3$, hence codimension $2$ in the fivefold
$F\times Y$.

\begin{proposition}\label{prop:Mprops}
$M$ is a simply connected compact K\"ahler manifold of complex dimension
$5$ with a Hamiltonian circle action lifting $\id_F\times(\text{diagonal
action})$. Its half-turn $\hat g$ is a Hamiltonian involution with
\[
  \Fix(\hat g)=C_1'\sqcup C_2'\sqcup C_3'\sqcup C_4',\qquad
  C_j'\cong\Bl_{S\times\{\mathrm{pt}\}}(F\times\CP^1),
\]
each $C_j'$ simply connected of real dimension $8$. The minimal Chern number
of $M$ is $1$.
\end{proposition}

\begin{proof}
K\"ahlerness and simple connectivity follow from
Proposition~\ref{prop:blowup}. $W$ is invariant, so the holomorphic circle
action lifts; averaging a K\"ahler form and using $H^1(M;\R)=0$ makes it
Hamiltonian, and its half-turn $\hat g$ is a Hamiltonian involution lifting
$g=\id_F\times g_Y$. We have $\Fix(g)=\bigsqcup_jF\times E_j$ and
$W\cap\Fix(g)=S\times\{\text{poles}\}$. At $p=(s,y)$ with $y$ a pole,
$E^{\ne}_p=\{0\}\oplus T_yL\subset T_sS\oplus T_yL=T_pW$, so
\eqref{eq:transverse} holds. Lemma~\ref{lem:transverse} gives
$\Fix(\hat g)=\bigsqcup\widehat{F\times E_j}$ with
$\widehat{F\times E_j}\cong\Bl_{S\times\{y_j\}}(F\times E_j)$, $y_j$ the
unique pole on $E_j$. These are simply connected by
Proposition~\ref{prop:blowup}. Finally, for $x\in F\setminus S$ the sphere
$\{x\}\times E_{00}$ is disjoint from $W$ and has $c_1=1$.
\end{proof}

\subsection{Integral homology and torsion profiles}

By Proposition~\ref{prop:blowup} (codimension $2$ in both cases),
\begin{equation}\label{eq:HMC}
  H_*(M)\cong H_*(F\times Y)\oplus2\,H_*(S\times\CP^1)[2],\qquad
  H_*(C_j')\cong H_*(F\times\CP^1)\oplus H_*(S)[2].
\end{equation}
All Künneth factors other than $F,S$ have free homology, so no Tor terms
arise. Write $a=b_2(F)$ and $\beta=b_2(S)$.

\begin{table}[h]
\centering\small
\begin{tabular}{@{}c|cccccc|cccc@{}}
\toprule
 & \multicolumn{6}{c|}{$M$} & \multicolumn{4}{c}{$C_j'$}\\
degree & 2 & 3 & 4 & 5 & 6 & 7 & 2 & 3 & 4 & 5\\
\midrule
$d_2$ & 1 & 1 & 6 & 6 & 1 & 1 & 1 & 1 & 1 & 1\\
$d_5$ & 0 & 2 & 2 & 2 & 2 & 0 & 0 & 1 & 1 & 0\\
$t_i=\max$ & 1 & 2 & 6 & 6 & 2 & 1 & 1 & 1 & 1 & 1\\
\bottomrule
\end{tabular}
\caption{Primary torsion multiplicities $d_p(\Tor H_i)$. All torsion is
elementary abelian. $\sum_it_i(M)=18$, $\sum_it_i(C_j')=4$.}
\label{tab:torsion}
\end{table}

\begin{table}[h]
\centering\small
\begin{tabular}{@{}lccccccc@{}}
\toprule
degree & 0 & 2 & 4 & 6 & 8 & 10 & total\\
\midrule
$r_i(M)$ & 1 & $a+8$ & $7a+2\beta+3$ & $7a+2\beta+3$ & $a+8$ & 1 & $16a+4\beta+24$\\
$r_i(C_j')$ & 1 & $a+2$ & $2a+\beta$ & $a+2$ & 1 & & $4a+\beta+6$\\
\bottomrule
\end{tabular}
\caption{Free ranks (all odd ranks vanish). $b(M)=4b(C_j')$.}
\label{tab:free}
\end{table}

\begin{proof}[Derivation of the tables]
$F$ has $\Tor H_2=\Tor H_3=\Z/2$; Künneth with $Y$ (free ranks $1,6,1$ in
degrees $0,2,4$) places $2$-torsion of multiplicities $(1,6,1)$ in degrees
$(2,4,6)$ and $(3,5,7)$. $S\times\CP^1$ has $5$-torsion $\Z/5$ in degrees
$1,2,3,4$; shifted by $2$ and doubled, this gives $(2,2,2,2)$ in degrees
$3,\dots,6$. For $C_j'$: $F\times\CP^1$ has $2$-torsion in degrees $2,3,4,5$,
and $H(S)[2]$ adds $\Z/5$ in degrees $3,4$. The free ranks follow in the
same way from $r(F)=(1,a,a,1)$, $r(S)=(1,\beta,1)$.
\end{proof}

\begin{corollary}\label{cor:counts}
\[
  \lambda(M)=b(M)+36,\qquad \lambda(C_j')=\tfrac14b(M)+8,\qquad
  \lambda(M)-\sum_{j=1}^4\lambda(C_j')=4,
\]
independently of $a$ and $\beta$. Moreover $\lambda_1(M)=b(M)+32=
\sum_j\lambda(C_j')$.
\end{corollary}

\begin{proof}
The first three identities follow from the tables. For $\lambda_1$, the
even-degree torsion of $M$ has $2$-multiplicity $1+6+1=8$ and
$5$-multiplicity $2+2=4$, so $d=8$; likewise in odd degrees. Thus
$\lambda_1(M)=b(M)+2(8+8)$.
\end{proof}

The source of the gap is visible in Table~\ref{tab:torsion}: in $M$ the
dominant prime is $5$ in degrees $3,6$ and $2$ in degrees $4,5$, so both
primes contribute; in $C_j'$ the $2$- and $5$-torsion occupy the same
degrees $3,4$ with equal multiplicity and are absorbed by the maximum.
The script \texttt{computations/homology/m10\_indep.py} reproduces these
numbers. It hard-codes placeholder free ranks $b(F)=8$, $b(S)=11$ and
prints $\lambda(M)=144$, $\sum\lambda(C_j')=140$, cyclic bound $140$; by
Corollary~\ref{cor:counts} the gap $4$ is independent of these
placeholders (also tested for $b(F)\in\{8,22\}$, $b_2(S)\in\{9,30\}$).
The actual values are given in Remark~\ref{rem:absolute}.

\begin{remark}\label{rem:absolute}
With $a=19$ (Proposition~\ref{prop:Fhom}) and $\beta=40$
(Proposition~\ref{prop:S}): $b(M)=488$, $\lambda(M)=524$,
$\lambda(C_j')=130$, $\#\Fix=520$, and the relative deficit is
$4/524=1/131$. The value $b(F)=40$ was obtained twice, from two
independent constructions of the toric fan (seven rays, nine smooth
maximal cones; \texttt{computations/sage/bF2.sage}); it depends on the
choice of resolution, the gap does not. Over $\F_2$ one has
$b(M;\F_2)=b(M)+2\cdot16=\#\Fix$, so the count also equals the total
mod-$2$ Betti number; over $\F_5$, $b(M;\F_5)=b(M)+16<\#\Fix$.
\end{remark}

\subsection{Proof of Theorem~\ref{thmA}}

By Theorem~\ref{thm:smale} with $D=8$, each $C_j'$ carries a Morse
function $h_j$ with exactly $\lambda(C_j')$ critical points.
Proposition~\ref{prop:perturb} applied to the
Hamiltonian involution $\hat g$ of Proposition~\ref{prop:Mprops} yields a
one-periodic Hamiltonian $H$ on $M$ with
$\#\Fix(\phi_H^1)=\sum_j\lambda(C_j')$, all nondegenerate, all with
contractible orbits since $M$ is simply connected. By
Corollary~\ref{cor:counts} this equals $\lambda(M)-4$, and
$\lambda(M)=\SM(M)=\Morse(M)$ by Theorem~\ref{thm:smale} with $D=10\ge9$
(\cite[Thm.~2.4]{OpenAI22}). The minimal Chern number is $1$
(Proposition~\ref{prop:Mprops}) and $\#\Fix=\lambda_1(M)$ by
Corollary~\ref{cor:counts}. \qed

\begin{remark}
Theorem~\ref{thmA} is consistent with Theorem~\ref{thm:BX}, which here is
an equality. It shows that the cyclic integral Floer bound is sharp in
complex dimension $5$ (previously complex dimension $1666$
\cite{OpenAICompanion}).
\end{remark}

%======================================================================
\section{Dimension eight: necessary conditions and a mechanism obstruction}\label{sec:eight}

\subsection{Necessary torsion conditions}

\begin{proposition}\label{prop:nec8}
Let $M$ be a simply connected closed symplectic $8$-manifold with minimal
Chern number $N$, and put $T_i=\Tor H_i(M)$. Then $T_i\cong T_{7-i}$,
$T_i=0$ for $i\notin\{2,3,4,5\}$, $\lambda_N(M)=\lambda(M)$ for $N\ge2$
(including $N=\infty$), and
\[
  \lambda(M)-\lambda_1(M)=4\bigl[d(T_2)+d(T_3)-d(T_2\oplus T_3)\bigr].
\]
Consequently, assuming Theorem~\ref{thm:BX} and Theorem~\ref{thm:smale}
for $D=8$, a nondegenerate Hamiltonian on $M$ with fewer than $\SM(M)$
fixed points can exist only if $N=1$ and $T_2,T_3$ are nonzero with
different dominant primes in the sense of \eqref{eq:dsum}. In particular
$\Tor H^3(M)\cong T_2\ne0$.
\end{proposition}

\begin{proof}
$T_i\cong\Tor H^{i+1}\cong T_{7-i}$ by universal coefficients and duality;
$T_0=T_1=0$ forces $T_6=T_7=0$. If $2N\ge4$ the degrees $2,3,4,5$ lie in
distinct classes mod $2N$. For $N=1$ the even class carries
$T_2\oplus T_4\cong T_2\oplus T_3$ and the odd class
$T_3\oplus T_5\cong T_3\oplus T_2$; compare with $\lambda$, which counts
$d(T_2)+d(T_3)$ twice in each parity. The rest follows from
Theorems~\ref{thm:BX} and~\ref{thm:smale} and \eqref{eq:dsum}.
\end{proof}

\begin{remark}
Blowing up a space with free homology along any centers never creates
torsion in degree $2$ (the shifts in Proposition~\ref{prop:blowup} are
$\ge2$ and $H_0$ of a center is free). Heuristically, therefore, a
blow-up construction in dimension $8$ needs an ambient with Brauer-type
torsion from the start; this remark applies only to constructions by
blow-ups.
\end{remark}

\begin{proposition}[failure of the product mechanism below dimension $12$]\label{prop:productfail}
Let $X$ be a simply connected closed K\"ahler manifold of real dimension
$\le6$ and $Y$ as in \S\ref{sec:ten}. Then
$\lambda(X\times Y)=4\lambda(X\times\CP^1)$.
\end{proposition}

\begin{proof}
If $\dim_\R X\le4$, $X$ has no torsion. If $\dim_\R X=6$, the proof of
Proposition~\ref{prop:six} shows that the only torsion is $T:=T_2\cong T_3$. Then
$X\times\CP^1$ has torsion $T$ in degrees $2,3,4,5$ and $X\times Y$ has
$T,T,T^6,T^6,T,T$ in degrees $2,\dots,7$, so
$\sum t_i(X\times Y)=16d(T)=4\sum t_i(X\times\CP^1)$; free ranks scale by
$4$ as well.
\end{proof}

\subsection{Perfectness of moment maps over prime fields}

\begin{lemma}\label{lem:perfect}
Let $C$ be a closed symplectic manifold with a Hamiltonian circle action,
moment map $\mu$, and fixed components $F$ of Morse--Bott index $2\mu_F$.
For every field $k$ and every $i$,
\[
  \dim H_i(C;k)=\sum_F\dim H_{i-2\mu_F}(F;k).
\]
\end{lemma}

\begin{proof}
Let $q$ be a prime. ``$\le$'': the Morse--Bott inequalities for $\mu$ hold
with $\F_q$ coefficients, since negative normal bundles are complex,
hence oriented, so the Thom isomorphism applies. ``$\ge$'' for total
dimensions: let $X_n=C^{\Z/q^n}$, where $\Z/q^n\subset S^1$. The group
$\Z/q^{n+1}$ preserves $X_n$ and acts on it through
$\Z/q^{n+1}/\Z/q^n\cong\Z/q$ with fixed set $X_{n+1}$. The Smith
inequality for $\Z/q$ acting smoothly on the closed manifold $X_n$
\cite[Prop.~III.4.16]{tomDieck1987} (stated for smooth $(\Z/q)^r$-manifolds
as \cite[Prop.~2.1]{Franz2016}; see also \cite[Ch.~III]{Bredon1972}) gives
$\dim H_*(X_{n+1};\F_q)\le\dim H_*(X_n;\F_q)$. A smooth action of $S^1$ on a compact manifold has
finitely many isotropy groups; since $\bigcup_n\Z/q^n$ is dense in $S^1$,
$X_n=C^{S^1}$ for $n$ large. Hence
$\sum_F\dim H_*(F;\F_q)\le\dim H_*(C;\F_q)$. Equality of total dimensions
forces the Morse--Bott spectral sequence (with $E^1$ the shifted homology of
the $F$'s) to degenerate, giving degreewise equality. For $k$ of
characteristic $0$, apply the result for a prime $q$ that divides no
torsion order of $C$ or of any $F$ (or use \cite{Frankel1959,Kirwan1984}).
The case of a general field reduces to its prime field.
\end{proof}

Lemma~\ref{lem:perfect} combines the classical Smith inequality (see
also \cite[Ch.~3]{AlldayPuppe1993}, \cite[Ch.~III]{Hsiang1975} for torus
actions) with the Morse--Bott inequalities, and we do not claim novelty
for it. We could not, however, find an explicit statement in the
literature that moment maps of Hamiltonian circle actions are perfect over
every $\F_p$. Perfectness is usually stated over $\Q$ or $\R$
\cite{Frankel1959,Kirwan1984}, or over $\Z$ when the fixed points are
isolated. The Atiyah--Bott--Kirwan argument via equivariant Euler classes
does not apply over $\F_p$ when $p$ divides an isotropy weight, which is
why we argue via Smith theory instead.

\begin{corollary}\label{cor:torsionformula}
In the situation of Lemma~\ref{lem:perfect}, for every prime $p$ and every
$i$,
\[
  r_i(C)=\sum_Fr_{i-2\mu_F}(F),\qquad
  t^p_i(C)=\sum_Ft^p_{i-2\mu_F}(F).
\]
\end{corollary}

\begin{proof}
$\dim H_i(\,\cdot\,;\F_p)=r_i+t^p_i+t^p_{i-1}$. Subtract the rational
identity and induct on $i$ starting from $t^p_{-1}=0$.
\end{proof}

\subsection{Proof of Theorem~\ref{thmB}}

Choose a circle $S^1\subset T$ with $M^{S^1}=M^T$: $T$ has finitely many
isotropy subgroups on $M$, and a circle whose Lie algebra avoids the
finitely many proper subspaces $\operatorname{Lie}(H)$, $H\ne T$ an
isotropy group, works. Let $\mu$ be its moment map. Each component $C$ of
$\Fix(g)$ is a $T$-invariant symplectic submanifold, $\mu|_C$ is a moment map
for the circle on $C$, and $C^{S^1}=C\cap M^T$. Write $\mathcal F$ for the
set of components of $M^T$ and $\mathcal F(C)\subset\mathcal F$ for those in
$C$; $\mathcal F=\bigsqcup_C\mathcal F(C)$. For $F\in\mathcal F(C)$ write
$2\nu_F$ and $2\mu_F$ for its indices in $C$ and in $M$.

\emph{Step 1: an upper bound for $\lambda(M)$.} By
Corollary~\ref{cor:torsionformula} applied to $M$,
$b(M)=\sum_{F}b(F)$ and
\[
  t_i(M)=\max_p\sum_Ft^p_{i-2\mu_F}(F)\le\sum_F\max_pt^p_{i-2\mu_F}(F)
  =\sum_Ft_{i-2\mu_F}(F).
\]
Summing over $i$: $\lambda(M)\le\sum_{F\in\mathcal F}\lambda(F)$.

\emph{Step 2: equality for each $C$.} We claim
$\lambda(C)=\sum_{F\in\mathcal F(C)}\lambda(F)$ for every component $C$ of
$\Fix(g)$; by hypothesis $\dim_\R C\le6$. (If $\dim_\R M=8$, this holds
unless $g$ acts trivially on $M$, in which case $C=M$ and the theorem is
trivial.) Nothing below uses $\dim M$ or $\pi_1(C)$. A component $F$
with nonzero torsion has $\dim_\R F\ge4$ (closed orientable manifolds of
dimension $\le2$ have free homology). If $F=C$ there is nothing to prove.
Otherwise $\dim_\R C=6$ and $\dim_\R F=4$, and the normal bundle of $F$ in
$C$ is a complex line, so $\nu_F\in\{0,1\}$ and $F$ is a local minimum or
local maximum of $\mu|_C$. By the connectedness theorem of Atiyah and
Guillemin--Sternberg \cite{Atiyah1982,GuilleminSternberg1982}, local
extrema are global and the minimum and maximum components are unique.
Hence $C$ contains at most two torsion-carrying fixed components: the
minimum $F_-$ ($\nu=0$, torsion in degrees $1,2$) and the maximum $F_+$
($\nu=1$, torsion in degrees $3,4$). These degree supports are disjoint,
so by Corollary~\ref{cor:torsionformula} each $t_i(C)$ equals $t_i(F_-)$ or
$t_{i-2}(F_+)$ or $0$, whence $\sum_it_i(C)=\sum_{F\in\mathcal F(C)}\sum_i
t_i(F)$. Together with $b(C)=\sum_{F\in\mathcal F(C)}b(F)$ this proves the
claim.

\emph{Step 3.} $\sum_j\lambda(C_j)=\sum_j\sum_{F\in\mathcal F(C_j)}\lambda(F)
=\sum_{F\in\mathcal F}\lambda(F)\ge\lambda(M)$.

\emph{Consequence.} The fixed points of $\psi^\epsilon\circ g$ are the
critical points of Morse functions $h_j$ on $C_j$, and
$\#\Crit(h_j)\ge\lambda(C_j)$ by Proposition~\ref{prop:SMlower}. Hence
$\#\Fix\ge\lambda(M)$, which equals $\SM(M)$ for simply connected $M$ by
Theorem~\ref{thm:smale} ($D=8$). \qed

\begin{remark}[a support criterion]\label{rem:support}
The proof shows, in any dimension, that $\sum_j\lambda(C_j)\ge\lambda(M)$
whenever, in every $C_j$, the index-shifted torsion supports of the
torsion-carrying fixed components are pairwise disjoint. A construction of
this type in dimension $2n$ therefore needs a component $C_j$ with two
torsion-carrying $T$-fixed components whose shifted supports overlap.
In Theorem~\ref{thmA} this happens as follows. The diagonal circle
has weights $(1,1)$ at $(0,0)$ and $(-1,-1)$ at $(\infty,\infty)$, so it
acts \emph{trivially} on $E_{00}$ and $E_{\infty\infty}$; hence $C_1'$ and
$C_4'$ (over these curves) are themselves $T$-fixed components. On $C_2'$
and $C_3'$ (over $E_{0\infty}$, $E_{\infty0}$) the circle rotates the
$\CP^1$-factor, and the $T$-fixed components are three: the proper
transform of $F\times\{y_j\}$ ($\cong F$, since $S\subset F$ is a divisor),
$F\times\{y_j'\}$ for the other fixed point $y_j'$ of $E_j$, and the
section $\PP(T_{y_j}E_j)\cong S$ of the exceptional divisor (the section
$\PP(N_{S/F})$ lies in the proper transform of $F\times\{y_j\}$). As a
check, the total Betti number of $M^T$ is
$2\cdot122+2\,(40+40+42)=488=b(M)$, as Lemma~\ref{lem:perfect} requires.
Within $C_2'$ one has $\sum_F\lambda(F)=44+44+46=134>130=\lambda(C_2')$:
the loss of $4$ comes from the overlap of the $5$-torsion of $S$ (degrees
$3,4$ after the index shift) with the $2$-torsion of the copies of $F$
(degrees $2$--$5$). Altogether $\sum_{F\subset M^T}\lambda(F)=528\ge
\lambda(M)=524\ge\sum_j\lambda(C_j')=520$, consistent with Step~1.
\end{remark}

\begin{remark}[a second argument via Li's theorem]
By a theorem of H.~Li \cite[Thm.~0.1]{Li2003}, for a Hamiltonian circle
action on a compact symplectic manifold $C$,
$\pi_1(C)\cong\pi_1(F_{\min})\cong\pi_1(F_{\max})\cong\pi_1(C_{\mathrm{red}})$. If $\dim C=6$ and $C$ is simply
connected, its four-dimensional fixed components (necessarily extremal)
are simply connected, hence torsion-free, and the torsion-carrying
fixed components of $M$ are themselves components of $\Fix(g)$; the
inequality of Step~1 then gives the claim. If some $C_j$ is not simply
connected, Lemma~\ref{lem:perfect} prevents cancellation of its
$H_2$-torsion. This gives an alternative route to Theorem~\ref{thmB}; we
have not written it out in full.
\end{remark}

\subsection{Dimensional limitation of transverse centers (heuristic)}\label{subsec:dimlimit}

The following is a heuristic, not a theorem. In Theorem~\ref{thmA} the extra torsion of the fixed components is
introduced by a blow-up whose center $W$ meets $\Fix(g)$ cleanly and
satisfies \eqref{eq:transverse}. If the inserted torsion is that of a
surface ($S$, real dimension $4$, with odd-prime torsion in $H_1$), then
$W$ must contain $S$ and, by \eqref{eq:transverse}, the normal directions
$E^{\ne}$ of the fixed component (real dimension $\ge2$), and must have
real codimension $\ge4$ (a divisor blow-up is the identity). Hence
$\dim_\R M\ge4+2+4=10$. The argument assumes that the inserted torsion
comes from a surface contained in the center; it is independent of, and
weaker than, Theorem~\ref{thmB}.

%======================================================================
\section{Higher dimensions, earlier claims, and open problems}\label{sec:higher}

\subsection{K\"ahler examples in dimensions twelve and sixteen}\label{subsec:higher}

\begin{remark}[dimension $12$]\label{rem:dim12}
Let $X^8=\Bl_{S\times\{t_1\}\sqcup S\times\{t_2\}}(F\times\CP^1)$ and
$M^{12}=X^8\times Y$, $C=X^8\times\CP^1$. Then $M^{12}$ is simply connected
K\"ahler, and the product mechanism gives
$\lambda(M^{12})-4\lambda(C)=8$; with $b(F)=40$, $b(S)=42$ one gets
$\lambda(M^{12})=1400$ and $\#\Fix=1392$. The gap relies on
Lemma~\ref{lem:torF} (no $5$-torsion in $F$); with $5$-torsion of
rank $\ge1$ in $F$ it would vanish.
\end{remark}

\begin{remark}[dimension $16$]\label{rem:dim16}
In $P=\CP^1\times\CP^5$ embed disjointly (by
\cite[Lemma~3.2]{OpenAI22}) a bielliptic surface $S_2$ with $(\Z/2)^2$
torsion and three bielliptic surfaces $S_3^{(i)}$ with $\Z/3$ torsion
into $\CP^5$, and blow up $P$ along $S_2\times\CP^1$ (codimension $3$) and
$S_3^{(i)}\times\{t\}$ (codimension $4$). For the resulting complex sixfold $X$ and
$M=X\times Y$ one finds $\lambda(M)=1232$, $4\lambda(X\times\CP^1)=1216$
(two independent scripts, \texttt{computations/homology/indep.py} and
\texttt{computations/homology/x12.py}); the
family with $3m$ copies of $S_3$ and $m$ copies of $S_2\times\CP^1$ has
$\lambda=96+1136m$ and gap $16m$. Here $\#\Fix=1216$ also equals the Bai--Xu cyclic bound
$\lambda_1(M)$ (minimal Chern number $1$), as recomputed with
\texttt{computations/homology/indep.py}. This is a rearrangement of the centers of
\cite{OpenAI22} and is not claimed as a separate result.
\end{remark}

\subsection{Relation to an earlier claim}\label{subsec:Ma}

For a closed symplectic manifold $(M,\omega)$, Ma \cite{Ma2013} writes
$\mathrm{Arn}(M,\omega)$ for the minimal number of fixed points of a
Hamiltonian symplectomorphism of $M$ and $\mathrm{Arn}_s(M,\omega)$ for
the minimum over those all of whose fixed points are nondegenerate. He
writes $\mathrm{Crit}(M)$ for the minimal number of critical points of a
smooth function and $\mathrm{Crit}_m(M)$ for the minimal number of
critical points of a Morse function. His Theorem~1.5 reads: ``Let
$(M,\omega)$ be a close[d] symplectic manifold with symplectic form
$\omega$. Then $\mathrm{Arn}(M,\omega)=\mathrm{Crit}(M)$ and
$\mathrm{Arn}_s(M,\omega)=\mathrm{Crit}_m(M)$'' \cite[Thm.~1.5]{Ma2013}.
There are no further hypotheses, and all fixed points are counted, not
only those with contractible orbits. See also \cite{VibeMathed}, cited for
context only.

The second equality is already refuted by \cite[Thm.~1.1]{OpenAISept23},
on a closed symplectic manifold of real dimension $12$ with $\pi_1=A_5^k$,
which is not simply connected, and by the simply connected examples of
\cite{OpenAIStable,OpenAI22,OpenAICompanion} in higher dimensions. The
first equality is refuted by the degenerate example of
\cite{OpenAIDegenerate} on the quadric $Q^3$. For the manifold $M$ of
Theorem~\ref{thmA}, $\mathrm{Crit}_m(M)=\Morse(M)=\SM(M)=\lambda(M)$ by
Theorem~\ref{thm:smale}, while $\#\Fix(\phi_H^1)=\lambda(M)-4$, so
$\mathrm{Arn}_s(M,\omega)\le\mathrm{Crit}_m(M)-4$. Theorem~\ref{thmA} is
thus the first simply connected example below the stable Morse number in
real dimension less than $22$, and it refutes the second assertion of
\cite[Thm.~1.5]{Ma2013} on a simply connected K\"ahler manifold of real
dimension $10$.

\subsection{Open problem: dimension eight}

By Theorems~\ref{thmA} and~\ref{thmB} and Proposition~\ref{prop:six}, among
simply connected closed symplectic manifolds of real dimension at least
$6$ the minimal real dimension of a counterexample to \eqref{eq:SMconj}
is $8$ or $10$ (Corollary~\ref{cor:mindim}). Simply connected real
dimension $4$ remains open in the sense of \S\ref{sec:six}, and
non-simply-connected manifolds of real dimension $4$, $6$, $8$, $10$ are
not addressed here.

\begin{problem}\label{prob:dim8}
Does there exist a closed (simply connected, K\"ahler) symplectic
$8$-manifold $M$ and a nondegenerate Hamiltonian diffeomorphism with fewer
than $\SM(M)$ fixed points?
\end{problem}

By Proposition~\ref{prop:nec8}, such an $M$ must have $N=1$ and
Brauer-type torsion $T_2\ne0$ whose dominant primes differ from those of
$T_3$. By Theorem~\ref{thmB}, the Hamiltonian cannot be obtained by
perturbing a finite-order element of a Hamiltonian torus action. A
positive answer requires either a finite-order Hamiltonian diffeomorphism
not contained in any Hamiltonian torus action, or an entirely different
mechanism. A negative answer requires a Floer-theoretic obstruction
stronger than Theorem~\ref{thm:BX}, which by Proposition~\ref{prop:nec8} is
not sufficient in general.

%======================================================================
\appendix
\section{Computational corroboration}\label{app:milp}

The script \texttt{computations/milp/lpsearch2.py} formulates a mixed-integer linear program
modelling the torsion bookkeeping of the finite-order/torus mechanism in
real dimension $2n$. Variables are nonnegative weights $x$ of
``torsion pieces'' $(\text{prime}\in\{p,q\},f,(a,b),c,j,\nu,\varepsilon)$: a
pair of Poincar\'e-dual cyclic summands in degrees $a<b$ of a $T$-fixed
component of complex dimension $f$, contained in a component $C$ of
$\Fix(g)$ of complex dimension $c$ with label $j\le K$, with index $2\nu$
in $C$ and $2(\nu+\varepsilon)$ in $M$. Constraints: Poincar\'e duality of
the torsion profiles of $M$ and of each $C$; no torsion of $M$ in degrees
$0,1$; uniqueness of minimum and maximum components in each $C$ and
globally; if $F=C$ then $C$ contains no other piece; total weight $\le1$.
The objective is the normalized gain
$\sum_k\max(\text{profile}_M(k))-\sum_C\sum_k\max(\text{profile}_C(k))$,
linearized with binary selectors. The modelling of the torsion of $C$ as
the sum of the shifted torsion of its fixed components is justified by
Corollary~\ref{cor:torsionformula}, so no extra perfectness hypothesis is
needed.

\begin{table}[h]
\centering\small
\begin{tabular}{@{}lll@{}}
\toprule
real dim & setting & max normalized gain\\
\midrule
$8$ & $K=2,3$; with or without $H_1$-torsion in $C$ & $0$\\
$10$ & $K=2,3$, no $H_1$-torsion in $C$ & $1/3$\\
$10$ & $K=3$ unrestricted & $1/2$\\
$12$ & $K=2,3$ & $1$ (positive control)\\
$10$ & forbid $6$-dim $F$ / Brauer-type $(2,3)$ $F$ / surfaces / $8$-dim $C$ & $0$ in each case\\
\bottomrule
\end{tabular}
\caption{Results (\texttt{computations/milp/run\_lp.py}, \texttt{computations/milp/necessity.py}).}
\end{table}

In dimension $10$ the optimal configurations contain an $8$-dimensional
$C$ whose minimum and maximum (or saddles) are $6$-dimensional
Brauer-type components and which contains torsion surfaces of the other
prime at index $2$; by Guillemin--Sternberg birational cobordism
\cite{GuilleminSternberg1989}, such a surface is a codimension-$2$
symplectic submanifold of a reduced space of Brauer type, which led to the
construction of \S\ref{sec:F}. \emph{Caveats:} this is an LP relaxation
with real weights and two primes, it encodes only the listed necessary
conditions and not realizability, and it is used only as corroboration;
Theorem~\ref{thmB} has a complete hand proof.

\section{Machine verification of the symmetroid and the quintic}\label{app:sage}

All paths are relative to the repository root. In exact arithmetic over
$\Q$ (and over $\overline{\Q}$ where stated):
\begin{itemize}
\item \texttt{computations/sage/q10.sage} (random seed $7$) produces the web
  and the quintic and verifies: the character $f(gx)=\zeta^3f(x)$;
  irreducibility of $f$ over $\Q$; that the saturated Jacobian ideal of $f$
  has Hilbert polynomial $10$ and equals its radical and the saturated
  ideal of $3\times3$ minors of $A$; that the coordinate points lie on $B$
  and are smooth points (gradients $(0,0,-1728,0)$, $(0,0,0,-875)$,
  $(0,-640,0,0)$, $(-672,0,0,0)$ at $p_0,\dots,p_3$); that the web has no
  base points; that a general (non-invariant) web also has total Tjurina
  number $10$; that the invariant quintic $Q$ is smooth, avoids the
  coordinate points and the nodes, and that $B\cap\{Q=0\}$ is a smooth curve
  (Hilbert polynomial $20t-50$: degree $20$, genus $51$).
\item \texttt{computations/sage/beauville\_ii.sage}: the rank-$\le1$ locus
  of the web is empty, and Beauville's condition (ii) holds at each of the
  ten rank-$2$ points, computed exactly over $\overline{\Q}$.
\item \texttt{computations/sage/defect2.sage}: the ideal of the nodes is
  radical of degree $10$, and its Hilbert function is $1,4,10,10,\dots$, so
  the nodes impose independent conditions on quadrics.
\item \texttt{computations/sage/bF2.sage}: the toric resolution of
  $\frac15(1,1,4)$ by iterated star subdivision has $7$ rays and $9$ smooth
  maximal cones; an independent construction of the fan gives the same
  result.
\item \texttt{computations/sage/symm.sage}: auxiliary experiments with
  random invariant webs.
\end{itemize}
The homology bookkeeping of \S\ref{sec:ten} is implemented twice, with
identical results: \texttt{computations/homology/lam.py} together with
\texttt{w10.py} tracks torsion profiles, and
\texttt{computations/homology/indep.py} together with
\texttt{m10\_indep.py} tracks degreewise integral groups with Künneth.
\texttt{computations/homology/x12.py} contains the higher-dimensional
scans. The MILP of Appendix~\ref{app:milp} is
\texttt{computations/milp/lpsearch2.py}, run by
\texttt{computations/milp/run\_lp.py} and \texttt{computations/milp/necessity.py}.

\section*{Statement on the use of AI}
The results in this manuscript were found, checked and written up by
Claude (Anthropic) agents working under the direction of the author.
The computations listed above were run by those agents. The manuscript
has not undergone human peer review.

\printbibliography

\end{document}
